\cotefile{156-4}
\scanpage{2}
\mspage{1}
\margin{\begin{tabular}{|l}\unsure{GF|CF} $\overline{\underline{\mathrm{IV}}}$\\\hline\end{tabular}}
\uline{Analysis situs} (\unsure{première mouture}).

(10 juin 1986)

1. Comme modèle ``\unsure{ensembliste} topologique de l'axiomatique d'un ``\unsure{modèle|monde} de \unsure{formes|forme} \inserted{\unsure{topologiques|topologique}}'' (ou d'une ``\uline{\unsure{contrée}}''),
considérons les notions \unsure{suivantes} \unsure{attachées} \unsure{à|:} un ``espace'' $\pm$ \unsure{ordinaire} $\deleted{\unsure{X}}\unsure{E|X}$.
On supposera ici que $\deleted{X}E$ \unsure{est} un espace topologique séparé, \unsure{ou} un espace modéré (pas néc. compact) pour une théorie \unsure{modérée} donnée.

Dans ce cas ``partie fermée'' signifie ``partie fermée \unsure{modérée}''

I) \uline{Figure} \qquad Une \uline{figure} \inserted{$F$} de $\deleted{X}E$ est une partie fermée $X$, \uline{munie} d'une
% slanted
\margin{$X$ est appelé le \uline{support} de la figure $\operatorname{supp} F$ ou $\lvert F\rvert$}
décomposition de $Y$ en \deleted{\unsure{strates}} \inserted{parties} $X_i$
\[
\{X_{\unsure{i}}\}\in\unsure{\mathcal{P}|D}(Y) \qquad X=\bigcup X_i
\]
appelées les \uline{strates fermées}. On fait les hypothèses

a) Les $X_i$ fermés \inserted{non vides et loc. compacts (loc. \unsure{cpt|cpt.} est \unsure{autom.} si $\deleted{X}E$ l'est, par ex. \unsure{Esp. mod.})} (pas néc. compacts \unsure{\}} \deleted{\illegible{1}} \unsure{ils le s\textsuperscript{t}} si $X$ l'est\unsure{)}

b) La famille des $X_i$ est loc. finie
\inserted{(donc par a) b) \unsure{d'où|d} $\Rightarrow$ réunion \unsure{de|des} str.~f. est fermée\unsure{)}}
\margin{NB on pose $i\leq j\iff X_i\unsure{\subset|\in}X_j$ \quad $i<j\iff X_i\unsure{\subsetneqq|\nsubseteq}X_j$}

c) Si \deleted{\unsure{$X_i=$}} \deleted{\unsure{$\partial X_i=$}} $X_i\subset X_j$, $X_i\ne X_j$, (i.e. $i<j$) alors $X_i$ \uline{rare} dans $X_j$

d) Posant $\partial X_i=\deleted{\unsure{\bigcup_{X_j\subsetneq X_i}X_j}}\bigcup_{j<i}X_j$ (donc $\partial X_i$
